Beam offset correction refers to the re-calibration of $\alpha$ using the data from zero-field runs. Subsequently, re-calibration of $\alpha$ leads to the re-calibration of $p_T$. 
Only beam offset correction for Run15 \pAu was uploaded to MasterRecalibrator. 
Values for beam offset correction for Run14 \HeAu and Run15 \pAl were taken from \href{https://phenix-intra.sdcc.bnl.gov/WWW/p/forms/info/scripts/pview.cgi}{AN1509} and \href{https://phenix-intra.sdcc.bnl.gov/phenix/WWW/p/draft/dlario/Charged_Hadrons_pAl/AN_Charged_hadrons_in_pAl_v1.pdf}{1453} respectively. The correction is performed by subtracting the fit of zero-field runs $\Delta\alpha$:

\begin{equation}
   \Delta\alpha = \frac{dx}{R_{\text{DC}}}\sin{\varphi_{\text{DC}}} + \frac{dy}{R_{\text{DC}}}\cos{\varphi_{\text{DC}}}
\end{equation}

Where $R_{\text{DC}}$ is azimuthal angle of a track at a DC reference radius $R_{\text{DC}}$, $dx$, $dy$ - approximation parameters. 
Parameters $dx$ and $dy$ are presented in Table \ref{table:beam_offset_parameters} for different datasets. 
Since $p_T \sim 1/\alpha$, and having $p_T'$ to be uncorrected value, the transverse momentum was corrected as follows:

\begin{equation}
   p_T = \frac{\alpha}{\alpha - \Delta\alpha}p_T'
\end{equation}

\begin{table}[h!]
	\centering
	\begin{tabular}{c||c|c|c|c}
      \multirow{2}{*}{Dataset} & \multicolumn{2}{|c|}{East arm} & \multicolumn{2}{|c}{West arm} \\
      \cline{2-5} & dx & dy & dx & dy \\
      \hline
      \hline
      Run14 \HeAu, $\text{run}<416379$ & -0.175 & -0.036 & 0.173 & 0.050 \\
      Run14 \HeAu, $416379<\text{run}<416442$ & -0.180 & -0.043 & -0.015 & 0.039 \\
      Run14 \HeAu, $\text{run}>416442$ & -0.186 & -0.049 & 0.030 & 0.030 \\
      \hline
      Run15 \pAl & 0.070 & 0.100 & 0.162 & 0.061 \\
	\end{tabular}
   \caption{Beam offset correction parameters}
	\label{table:beam_offset_parameters}
\end{table}
